% ch06_numpy.tex -- NumPy: Arrays and Numerical Computing
\chapter{NumPy --- Arrays and Numerical Computing}
\label{ch:numpy}

\begin{intuition}
Imagine having to measure the temperature at 1,000 points on a hot plate.
With a Python list, each operation requires a loop. NumPy\index{NumPy}
lets you process \emph{all points in a single instruction}, like an instrument
that measures the entire plate at once. It is the difference between counting
grains of sand one by one and weighing the whole pile.
\end{intuition}

% ============================================================
\section{Why NumPy?}

\begin{codeblock}[title=\textbf{Python}]
\begin{minted}{python}
import time

# Comparison: sum of squares from 1 to 1,000,000
n = 1_000_000

# With a Python list
debut = time.time()
resultat_liste = sum([i**2 for i in range(n)])
temps_liste = time.time() - debut

# With NumPy
import numpy as np
debut = time.time()
resultat_numpy = np.sum(np.arange(n)**2)
temps_numpy = time.time() - debut

print(f"Python list: {temps_liste:.4f} s")
print(f"NumPy      : {temps_numpy:.4f} s")
print(f"Speedup    : x{temps_liste/temps_numpy:.0f}")
\end{minted}
\end{codeblock}

\begin{codeoutput}
\begin{verbatim}
Python list: 0.2851 s
NumPy      : 0.0032 s
Speedup    : x89
\end{verbatim}
\end{codeoutput}

% ============================================================
\section{Creating arrays}
\index{np.array}\index{np.zeros}\index{np.ones}\index{np.linspace}\index{np.arange}

\begin{codeblock}[title=\textbf{Python}]
\begin{minted}{python}
import numpy as np

# From a list
a = np.array([1.0, 2.0, 3.0, 4.0])
print(f"Array     : {a}")

# Special arrays
zeros = np.zeros(5)
print(f"Zeros     : {zeros}")

uns = np.ones(4)
print(f"Ones      : {uns}")

# Regular sequence
x = np.arange(0, 10, 2)
print(f"Arange    : {x}")

# N equally spaced points
t = np.linspace(0, 1, 5)
print(f"Linspace  : {t}")
\end{minted}
\end{codeblock}

\begin{codeoutput}
\begin{verbatim}
Array     : [1. 2. 3. 4.]
Zeros     : [0. 0. 0. 0. 0.]
Ones      : [1. 1. 1. 1.]
Arange    : [0 2 4 6 8]
Linspace  : [0.   0.25 0.5  0.75 1.  ]
\end{verbatim}
\end{codeoutput}

% --- TikZ array indexing ---
\begin{figure}[ht]
\centering
\begin{tikzpicture}[
    cell/.style={rectangle, draw, minimum width=1.2cm, minimum height=0.8cm,
        fill=blue!10, align=center},
    idx/.style={font=\small\ttfamily, text=red!70!black},
]
\foreach \val [count=\i from 0] in {1.0, 2.0, 3.0, 4.0, 5.0} {
    \node[cell] (c\i) at (\i*1.3, 0) {\val};
    \node[idx, above=2pt of c\i] {\i};
    \pgfmathtruncatemacro{\neg}{\i - 5}
    \node[idx, below=2pt of c\i] {\neg};
}
\node[font=\small, left=0.5cm of c0, text=red!70!black] {indices};
\end{tikzpicture}
\caption{Indexing a NumPy array: positive indices (top) and negative indices (bottom).}
\label{fig:numpy-indexation}
\end{figure}

% ============================================================
\section{Element-wise operations}
\index{element-wise operation}\index{broadcasting}

\begin{codeblock}[title=\textbf{Python}]
\begin{minted}{python}
import numpy as np

a = np.array([1, 2, 3, 4])
b = np.array([10, 20, 30, 40])

print(f"a + b   = {a + b}")
print(f"a * b   = {a * b}")
print(f"a ** 2  = {a ** 2}")
print(f"a + 100 = {a + 100}")    # broadcasting
print(f"a > 2   = {a > 2}")       # comparison
\end{minted}
\end{codeblock}

\begin{codeoutput}
\begin{verbatim}
a + b   = [11 22 33 44]
a * b   = [ 10  40  90 160]
a ** 2  = [ 1  4  9 16]
a + 100 = [101 102 103 104]
a > 2   = [False False  True  True]
\end{verbatim}
\end{codeoutput}

\begin{definition}[Broadcasting]
\textbf{Broadcasting}\index{broadcasting} allows NumPy to apply operations
between arrays of different sizes. A scalar is automatically
``extended'' to match the size of the array.
\end{definition}

% ============================================================
\section{Indexing and slicing}

\begin{codeblock}[title=\textbf{Python}]
\begin{minted}{python}
import numpy as np

x = np.array([10, 20, 30, 40, 50, 60, 70])

print(f"Element 3     : {x[3]}")
print(f"Slice [2:5]   : {x[2:5]}")

# Boolean indexing (fancy indexing)
masque = x > 35
print(f"Mask          : {masque}")
print(f"Values > 35   : {x[masque]}")

# Conditional modification
x[x < 30] = 0
print(f"After modif   : {x}")
\end{minted}
\end{codeblock}

\begin{codeoutput}
\begin{verbatim}
Element 3     : 40
Slice [2:5]   : [30 40 50]
Mask          : [False False False  True  True  True  True]
Values > 35   : [40 50 60 70]
After modif   : [ 0  0 30 40 50 60 70]
\end{verbatim}
\end{codeoutput}

% ============================================================
\section{Mathematical functions}
\index{np.sin}\index{np.cos}\index{np.exp}\index{np.log}\index{np.sqrt}

\begin{codeblock}[title=\textbf{Python}]
\begin{minted}{python}
import numpy as np

x = np.linspace(0, 2 * np.pi, 5)

print(f"x       = {np.round(x, 2)}")
print(f"sin(x)  = {np.round(np.sin(x), 4)}")
print(f"cos(x)  = {np.round(np.cos(x), 4)}")
print(f"exp(x)  = {np.round(np.exp(x), 2)}")

# Application: damped signal
t = np.linspace(0, 5, 6)
signal = np.exp(-0.5 * t) * np.sin(2 * np.pi * t)
print(f"\nt       = {np.round(t, 1)}")
print(f"signal  = {np.round(signal, 4)}")
\end{minted}
\end{codeblock}

\begin{codeoutput}
\begin{verbatim}
x       = [0.   1.57 3.14 4.71 6.28]
sin(x)  = [ 0.      1.      0.     -1.     -0.    ]
cos(x)  = [ 1.      0.     -1.     -0.      1.    ]
exp(x)  = [  1.     4.81  23.14 111.32 535.49]

t       = [0. 1. 2. 3. 4. 5.]
signal  = [ 0.     -0.      0.     -0.      0.     -0.    ]
\end{verbatim}
\end{codeoutput}

% ============================================================
\section{Descriptive statistics}
\index{np.mean}\index{np.std}\index{np.min}\index{np.max}

\begin{codeblock}[title=\textbf{Python}]
\begin{minted}{python}
import numpy as np

# Experimental measurements (concentration in mg/L)
mesures = np.array([12.3, 11.8, 12.5, 11.9, 12.1, 12.4, 12.0, 11.7])

print(f"Mean        : {np.mean(mesures):.2f} mg/L")
print(f"Std dev     : {np.std(mesures):.2f} mg/L")
print(f"Minimum     : {np.min(mesures):.2f} mg/L")
print(f"Maximum     : {np.max(mesures):.2f} mg/L")
print(f"Median      : {np.median(mesures):.2f} mg/L")
\end{minted}
\end{codeblock}

\begin{codeoutput}
\begin{verbatim}
Mean        : 12.09 mg/L
Std dev     : 0.27 mg/L
Minimum     : 11.70 mg/L
Maximum     : 12.50 mg/L
Median      : 12.05 mg/L
\end{verbatim}
\end{codeoutput}

% ============================================================
\section{Elementary linear algebra}
\index{np.dot}\index{matrix multiplication}

\begin{codeblock}[title=\textbf{Python}]
\begin{minted}{python}
import numpy as np

# 2x2 matrices
A = np.array([[1, 2],
              [3, 4]])

B = np.array([[5, 6],
              [7, 8]])

# Matrix product
C = A @ B    # equivalent to np.dot(A, B)
print("A @ B =")
print(C)

# Determinant and inverse
det_A = np.linalg.det(A)
inv_A = np.linalg.inv(A)
print(f"\ndet(A) = {det_A:.1f}")
print(f"A^(-1) =\n{inv_A}")
\end{minted}
\end{codeblock}

\begin{codeoutput}
\begin{verbatim}
A @ B =
[[19 22]
 [43 50]]

det(A) = -2.0
A^(-1) =
[[-2.   1. ]
 [ 1.5 -0.5]]
\end{verbatim}
\end{codeoutput}

% ============================================================
\section{Random numbers}
\index{np.random}

\begin{codeblock}[title=\textbf{Python}]
\begin{minted}{python}
import numpy as np

rng = np.random.default_rng(seed=42)

# Uniform between 0 and 1
u = rng.uniform(0, 1, size=5)
print(f"Uniform   : {np.round(u, 3)}")

# Normal distribution (mean=0, std=1)
n = rng.normal(0, 1, size=5)
print(f"Normal    : {np.round(n, 3)}")

# Random integers (dice roll)
des = rng.integers(1, 7, size=10)
print(f"Rolls     : {des}")
\end{minted}
\end{codeblock}

\begin{codeoutput}
\begin{verbatim}
Uniform   : [0.773 0.438 0.859 0.697 0.094]
Normal    : [ 0.314  0.459 -0.537  0.268 -0.232]
Rolls     : [5 2 3 6 1 4 2 5 3 6]
\end{verbatim}
\end{codeoutput}

% ============================================================
\section{Common errors}

\begin{errmark}
\begin{enumerate}
    \item \textbf{Dimension mismatch (\py{shape mismatch})}: operating on
    arrays of incompatible sizes causes a \py{ValueError}.
    \begin{minted}{python}
a = np.array([1, 2, 3])
b = np.array([1, 2])
# a + b -> ValueError: operands could not be broadcast
    \end{minted}

    \item \textbf{Forgetting element-wise operations}: the \py{*} operator
    between two arrays performs element-wise multiplication, \emph{not}
    matrix multiplication. Use \py{@} or \py{np.dot} for the matrix product.

    \item \textbf{Confusing \py{np.arange} / \py{np.linspace}}:
    \py{np.arange(0, 1, 0.1)} specifies a \emph{step},
    \py{np.linspace(0, 1, 10)} specifies a \emph{number of points}.

    \item \textbf{Modifying a view instead of a copy}: NumPy slicing creates a
    \emph{view} (not a copy). Modifying the view modifies the original!
    \begin{minted}{python}
a = np.array([1, 2, 3, 4])
b = a[1:3]     # view, not copy
b[0] = 99      # a becomes [1, 99, 3, 4]!
b = a[1:3].copy()  # independent copy
    \end{minted}
\end{enumerate}
\end{errmark}

% ============================================================
\section{Mini-project: Estimating $\pi$ with Monte Carlo}

\begin{projet}
The Monte Carlo method\index{Monte Carlo} estimates $\pi$ by drawing random points
in a square $[-1, 1]^2$ and counting those that fall inside the
unit circle ($x^2 + y^2 \leq 1$). The ratio gives:
\[
    \pi \approx 4 \times \frac{\text{points inside the circle}}{\text{total points}}
\]
\end{projet}

\begin{codeblock}[title=\textbf{Python}]
\begin{minted}{python}
import numpy as np

rng = np.random.default_rng(seed=42)
N = 1_000_000

# Draw N random points in [-1, 1] x [-1, 1]
x = rng.uniform(-1, 1, size=N)
y = rng.uniform(-1, 1, size=N)

# Count points inside the unit circle
dans_cercle = x**2 + y**2 <= 1
nb_dans_cercle = np.sum(dans_cercle)

# Estimate pi
pi_estime = 4 * nb_dans_cercle / N
erreur = abs(pi_estime - np.pi)

print(f"Total points       : {N:,}")
print(f"Points in circle   : {nb_dans_cercle:,}")
print(f"Estimated pi       : {pi_estime:.6f}")
print(f"Exact pi           : {np.pi:.6f}")
print(f"Error              : {erreur:.6f}")
\end{minted}
\end{codeblock}

\begin{codeoutput}
\begin{verbatim}
Total points       : 1,000,000
Points in circle   : 785,436
Estimated pi       : 3.141744
Exact pi           : 3.141593
Error              : 0.000151
\end{verbatim}
\end{codeoutput}

% ============================================================
\section{Exercises}

\begin{exercise}[$\star$ --- Array Creation]
Create the following arrays: (a) integers from 0 to 99, (b) 50 equally
spaced values between $-\pi$ and $\pi$, (c) a $3\times 3$ matrix of zeros
with 1s on the diagonal (identity matrix).
\end{exercise}

\begin{exercise}[$\star$ --- Temperature Conversion]
Create an array of temperatures in Celsius from $-40$ to $100$ (step of 10).
Convert it to Fahrenheit ($F = 1.8 \times C + 32$) \emph{without a loop}.
\end{exercise}

\begin{exercise}[$\star\star$ --- Experimental Data Analysis]
Generate 1000 simulated measurements from a normal distribution ($\mu = 50$, $\sigma = 5$).
Compute the mean, standard deviation, and count how many measurements are more than
$2\sigma$ from the mean. Compare with the theoretical prediction ($\approx 5\%$).
\end{exercise}

\begin{exercise}[$\star\star$ --- Dot Product]
Compute the angle between vectors $\vec{u} = (1, 2, 3)$ and $\vec{v} = (4, 5, 6)$
using the formula $\cos\theta = \frac{\vec{u}\cdot\vec{v}}{|\vec{u}||\vec{v}|}$.
\end{exercise}

\begin{exercise}[$\star\star\star$ --- Solving a Linear System]
Solve the following system of linear equations using \py{np.linalg.solve}:
\[
\begin{cases}
2x + 3y - z = 1 \\
4x + y + 2z = 2 \\
-x + 2y + 3z = 3
\end{cases}
\]
Verify your result by computing $A\vec{x}$ and comparing with $\vec{b}$.
\end{exercise}

% ============================================================
\begin{keyformulas}{Chapter Summary: NumPy}
\begin{itemize}
    \item \textbf{Creation}: \py{np.array}, \py{np.zeros}, \py{np.ones},
    \py{np.linspace}, \py{np.arange}
    \item \textbf{Operations}: \py{+}, \py{-}, \py{*}, \py{/}, \py{**}
    are applied element-wise
    \item \textbf{Broadcasting}: a scalar is automatically extended
    \item \textbf{Math}: \py{np.sin}, \py{np.cos}, \py{np.exp}, \py{np.log}, \py{np.sqrt}
    \item \textbf{Stats}: \py{np.mean}, \py{np.std}, \py{np.min}, \py{np.max}, \py{np.median}
    \item \textbf{Algebra}: \py{A @ B} (matrix product), \py{np.linalg.det},
    \py{np.linalg.inv}, \py{np.linalg.solve}
    \item \textbf{Random}: \py{rng = np.random.default\_rng(seed)}
    \item \textbf{Boolean indexing}: \py{x[x > 0]} filters positive elements
\end{itemize}
\end{keyformulas}
